
\documentclass[a4paper,10pt]{report}

%\usepackage[landscape]{geometry}
\usepackage[cm-default]{fontspec}
\setromanfont{FreeSerif}
\setsansfont{FreeSans}
\setmonofont{FreeMono}
\usepackage[utf8x]{inputenc}
\usepackage{fontspec}
\usepackage{xunicode}
\usepackage{xltxtra}
\usepackage{xgreek}
\usepackage{amsmath}
\usepackage{unicode-math}
\usepackage{ulem}
\usepackage{color}
\usepackage{verbatim}
\usepackage{nopageno}
\usepackage{graphicx}
\usepackage{textpos}
\setlength{\TPHorizModule}{1cm}
\setlength{\TPVertModule}{1cm}
%\usepackage[colorgrid,texcoord]{eso-pic}
\usepackage[outline]{contour}
\usepackage{wrapfig}
\usepackage{url}
\usepackage{color}
\usepackage{tikz}
\usepackage{fancybox,fancyhdr}
\usepackage{subfigure}
\usepackage{pstricks}
\usepackage{epsfig}
\usepackage{multicol}
\usepackage{listings}
\usepackage{enumerate}
\usepackage{hyperref}
\hypersetup{
  bookmarks=true,
  bookmarksopen=true,
  pdfborder=false,
  pdfpagemode=UseNone,
  raiselinks=true,
  pdfhighlight={/P},
  colorlinks,
  citecolor=black,
  filecolor=black,
  linkcolor=black,
  urlcolor=black
}


\usepackage{fontspec}
\usepackage{xunicode}
\usepackage{xltxtra}


% Margins
%\setlength{\textwidth}{24cm}

%\setlength{\voffset}{0in}
%\setlength{\textheight}{6.8in}
% Colors
\definecolor{cbrown}{rgb}{0.49,0.24,0.07}
\definecolor{cmpez}{rgb}{0.92,0.88,0.79}
\definecolor{cmpez2}{rgb}{0.62,0.33,0.34}
\definecolor{cwhite}{rgb}{1,1,1}
\definecolor{cblack}{rgb}{0,0,0}
\definecolor{cred}{rgb}{0.9,0.15,0.15}
\definecolor{clblue}{rgb}{0.57,0.85,0.97}
\definecolor{corange}{rgb}{0.86,0.49,0.18}
\definecolor{clorange}{rgb}{0.95,0.84,0.61}
\definecolor{cyellow}{rgb}{1,0.95,0.16}
\definecolor{cgreen}{rgb}{0.69,0.8,0.31}
\definecolor{clbrown}{rgb}{0.78,0.67,0.43}
\definecolor{cmagenta}{rgb}{0.79,0.05,0.54}
\definecolor{cgray}{rgb}{0.76,0.73,0.63}
\definecolor{bluesite}{rgb}{0.05,0.43,0.67}
\definecolor{redsite}{rgb}{0.90,0.15,0.15}
\definecolor{backsite}{rgb}{0.07,0.07,0.07}
% ----------- from Costas -----------------------------

%---------------   Main Fonts    -----------------%   

% \setmainfont[Mapping=tex-text]{Calibri}
% \setmainfont[Mapping=tex-text]{Times New Roman}
%\setmainfont[Mapping=tex-text]{Droid Serif}
%\setmainfont[Mapping=tex-text]{Cambria}
%\setmainfont[Mapping=tex-text]{cm-unicode}
% \setmainfont[Mapping=tex-text]{Gentium}
% \setmainfont[Mapping=tex-text]{GFS Didot}
% \setmainfont[Mapping=tex-text]{Comic Sans MS}
 \setmainfont[Mapping=tex-text]{Ubuntu}
% \setmainfont[Mapping=tex-text]{Myriad Pro}
% \setmainfont[Mapping=tex-text]{CMU Concrete}
%\setmainfont[Mapping=tex-text]{DejaVu Sans}
%\setmainfont[Mapping=tex-text]{KerkisSans}
%\setmainfont[Mapping=tex-text]{KerkisCaligraphic}
% \setmainfont[Mapping=tex-text]{Segoe Print}
% \setmainfont[Mapping=tex-text]{Gabriola}

%------------------------------------------------%

% ------------ Mathematics Fonts ----------------%
\setmathfont{Asana-Math.ttf}
% -----------------------------------------------%
% Misc
\author{Κασωτάκη Ε. - Σμαραγδάκης Κ.}
\title{www.math24.gr}
% ---------------- Formation --------------------%


\setlength\topmargin{-0.7cm}
\addtolength\topmargin{-\headheight}
\addtolength\topmargin{-\headsep}
\setlength\textheight{26cm}
\setlength\oddsidemargin{-0.54cm}
\setlength\evensidemargin{-0.54cm}
%\setlength\marginparwidth{1.5in}
\setlength\textwidth{17cm}

\RequirePackage[avantgarde]{quotchap}
\renewcommand\chapterheadstartvskip{\vspace*{0\baselineskip}}
\RequirePackage[calcwidth]{titlesec}
\titleformat{\section}[hang]{\bfseries}
{\Large\thesection}{12pt}{\Large}[{\titlerule[0.9pt]}]
%--------------------------------------------------%
\usepackage{draftwatermark}
\SetWatermarkText{www.math24.gr}
\SetWatermarkLightness{0.9}
\SetWatermarkScale{0.6}
%--------------------------------------------------%


\usepackage{xcolor}
\usepackage{amsthm}
\usepackage{framed}
\usepackage{parskip}

\colorlet{shadecolor}{bluesite!20}

\newtheorem*{orismos}{Ορισμός}

\newenvironment{notation}
  {\begin{shaded}\begin{theorem}}
  {\end{theorem}\end{shaded}}




% Document begins
\begin{document}
\pagestyle{fancy}
\fancyhead{}
\fancyfoot{}
\renewcommand{\headrulewidth}{0pt}
\renewcommand{\footrulewidth}{0pt}


\fancyhead[LO,LE]{
 %\textblockcolor{backsite}
 %\begin{textblock}{5}(-2,-0.55)
  %\rule{0cm}{1cm}
 %\end{textblock}
 \textblockcolor{bluesite}
 \begin{textblock}{5}(-1.5,-0.55)
  \rule{0cm}{1cm}
 \end{textblock}
 %\textblockcolor{bluesite}
 %\begin{textblock}{14}(5.5,-0.55)
 % \rule{0cm}{1cm}
 %\end{textblock}
 \begin{textblock}{0}(-1,-0.25)
 \color{cwhite} \begin{Large}www.math24.gr\end{Large}
 \end{textblock}
\begin{textblock}{0}(16,-0.4)
 \color{cwhite} \includegraphics[height=1.5cm]{math24_logo.png}
 \end{textblock}
\textblockcolor{white}
\begin{textblock}{17}(-1,28)
 \color{backsite} \begin{small}Copyright \textcopyright 2011, Κασωτάκη Ε.(ikasotaki@gmail.com) - Σμαραγδάκης Κ.(kesmarag@gmail.com)\end{small}
 \end{textblock}
}
 
\begin{shaded}
\begin{center}
\huge \textbf{Φυλλάδιο Ασκήσεων}\\
%Πρόσθεση ρητών αριθμών
\end{center} 
\textbf{Μαθηματικά Α' Γυμνασίου} \hfill \textbf{Ημερομηνία Παράδοσης : \hspace{2em} }
\subsection*{Ονοματεπώνυμο :\hfill  \hspace{5em}}
\end{shaded}
\vspace{2em}
\begin{itemize}
 \item Δεκαδικά κλάσματα 
 \item Σύγκριση δεκαδικών αριθμών
 \item Στρογγυλοποίηση δεκαδικών αριθμών 
\end{itemize}

\section*{Θεωρία - Δεκαδικά κλάσματα \hfill \small{}}
\textbf{Δεκαδικό κλάσμα} λέγεται το κλάσμα που έχει παρονομαστή μία δύναμη του 10.  \\
\textbf{π.χ} $\dfrac{7}{10}$\\
\textbf{π.χ} $\dfrac{9}{100}$\\
\textbf{π.χ} $\dfrac{81}{100}$\\



\section*{Άσκηση 1  \hfill \small{20 μονάδες}}
Να γράψετε ως κλάσματα τους παρακάτω δεκαδικούς αριθμούς:
\begin{enumerate}[1)]
\begin{multicols}{2}
 \item $7,49$
 \item $8,1$
 \item $0,003$
 \item $0,27$
 \item $0,1$
 \item $3,701$
 \item $0,009$
 \item $7,13$
 \item $5,1$
 \item $0,28$
\end{multicols}
\end{enumerate}

\section*{Άσκηση 2  \hfill \small{20 μονάδες}}
Να γράψετε τα παρακάτω κλάσματα ως δεκαδικούς αριθμούς:
\begin{enumerate}[1)]
\begin{multicols}{2}
 \item $\dfrac{7}{10}$
 \item $\dfrac{88}{10}$
 \item $\dfrac{95}{100}$
 \item $\dfrac{97}{1000}$
 \item $\dfrac{589}{10}$
 \item $\dfrac{43}{1000}$
 \item $\dfrac{5107}{10000}$
 \item $\dfrac{71}{1000}$
\end{multicols}
\end{enumerate}


\section*{Άσκηση 3  \hfill \small{20 μονάδες}}
Να γράψετε τα παρακάτω κλάσματα ως δεκαδικούς αριθμούς.\\
Υπόδειξη: Να εκτελέσετε τις αντίστοιχες διαιρέσεις.
\begin{enumerate}[i)]
 \item $\dfrac{40}{5}$
 \item $\dfrac{25}{40}$
 \item $\dfrac{142}{20}$
 \item $\dfrac{14}{4}$
 \item $\dfrac{27}{6}$
\end{enumerate}

\section*{Θεωρία - Σύγκριση Δεκαδικών Αριθμών\hfill \small{}}
\begin{itemize}
 \item Για να συγκρίνουμε δύο δεκαδικούς αριθμούς, ελέγχουμε την τάξη κάθε αριθμού. Μεγαλύτερος είναι 
       εκείνος που είναι μεγαλύτερης τάξης.\\
      \textbf{π.χ} $88,07 > 9,95$
 \item Αν οι δύο αριθμοί που θέλουμε να συγκρίνουμε είναι ίδιας τάξης, τότε μεγαλύτερος είναι εκείνος 
       που έχει το μεγαλύτερο ψηφίο στην αρχική τάξη.\\
       \textbf{π.χ} $73,02 > 63,07$
 \item Αν οι δύο αριθμοί που θέλουμε να συγκρίνουμε αρχίζουν από ίδιο ψηφίο ίδιας τάξης, τότε μεγαλύτερος 
       είναι εκείνος που έχει το αμέσως επόμενο ψηφίο μεγαλύτερο.\\
      \textbf{π.χ} $53,08 > 52,11$\\
      \textbf{π.χ} $483,71 > 483,49$
\end{itemize}



\section*{Άσκηση 4  \hfill \small{20 μονάδες}}
Να συγκρίνετε τους παρακάτω αριθμούς.
\begin{enumerate}[i)]
 \item $79,145$ και $89,148$
 \item $11,273$ και $11,412$
 \item $1,3782$ και $1,381$
 \item $5181,7$ και $687,19$
 \item $27,318$ και $27,331$
\end{enumerate}


\section*{Θεωρία - Στρογγυλοποίηση Δεκαδικών Αριθμών\hfill \small{}}
Για να στρογγυλοποιήσουμε ένα δεκαδικό αριθμό
\begin{itemize}
 \item \textbf{1ο Βήμα:} Προσδιορίζουμε τη δεκαδική τάξη στην οποία θα γίνει η στρογγυλοποίηση.
 \item \textbf{2ο Βήμα:} Εξετάζουμε το ψηφίο της αμέσως μικρότερης τάξης και διακρίνουμε περιπτώσεις: 
        \begin{itemize}
         \item \textbf{1η Περίπτωση:} Αν αυτό είναι μικρότερο του $5$, το ψηφίο αυτό και όλα τα ψηφία 
                των μικρότερων τάξεων μηδενίζονται.
         \item \textbf{2η Περίπτωση:} Αν είναι μεγαλύτερο ή ίσο του $5$, το ψηφίο αυτό και όλα τα ψηφία 
                των μικρότερων τάξεων μηδενίζονται και το ψηφίο της τάξης στρογγυλοποίησης αυξάνεται κατά $1$.
        \end{itemize}
\end{itemize}

\textbf{Παράδειγμα 1}\\
Αν θέλουμε να στρογγυλοποιήσουμε τον αριθμό $51,8397$ στο δέκατο τότε:
\begin{itemize}
 \item \textbf{1ο Βήμα:} Η στρογγυλοποίηση θα γίνει στο ψηφίο $8$ (το πρώτο ψηφίο μετά την υποδιαστολή)
 \item \textbf{2ο Βήμα:} Το ψηφίο της αμέσως μικρότερης τάξης είναι το $3$ (το δεύτερο ψηφίο μετά την υποδιαστολή) 
       και αφού το $3$ είναι μικρότερο από το $5$ τότε ο αριθμός στρογγυλοποιείται ως $51,8$
\end{itemize}

\textbf{Παράδειγμα 2}\\
Αν θέλουμε να στρογγυλοποιήσουμε τον αριθμό $2,4381$ στο εκατοστό τότε:
\begin{itemize}
 \item \textbf{1ο Βήμα:} Η στρογγυλοποίηση θα γίνει στο ψηφίο $3$ (το δεύτερο ψηφίο μετά την υποδιαστολή)
 \item \textbf{2ο Βήμα:} Το ψηφίο της αμέσως μικρότερης τάξης είναι το $8$ (το τρίτο ψηφίο μετά την υποδιαστολή) 
       και αφού το $8$ είναι μεγαλύτερο από το $5$ τότε ο αριθμός στρογγυλοποιείται ως $2,44$
\end{itemize}


\section*{Άσκηση 5  \hfill \small{20 μονάδες}}
Να στρογγυλοποιήσετε τους παρακάτω δεκαδικούς αριθμούς στο δέκατο και στο εκατοστό:
\begin{enumerate}[i)]
 \item $123,4567$
 \item $8,23408$
 \item $837,87651$
 \item $13,1313$
 \item $8173,47961$
\end{enumerate}




\end{document}
